% !TeX program = xelatex
% !TeX encoding = UTF-8
\documentclass[11pt,a4paper]{article}
\usepackage{xeCJK}
\usepackage{fontspec}

% Настройка шрифтов для японского языка
\setCJKmainfont{Kozuka Mincho Pro}
\setCJKsansfont{Kozuka Gothic Pro}
\setCJKmonofont{Kozuka Gothic Pro}

% Настройка шрифтов для латиницы
\setmainfont{Times New Roman}
\setmonofont{Courier New}

\usepackage{amsmath,amssymb,amsthm,mathtools}
\usepackage{geometry}
\usepackage{booktabs}
\usepackage{hyperref}
\usepackage{url}
\usepackage{tabularx}
\usepackage{array}
\usepackage{makecell}
\usepackage{ragged2e}
\usepackage{bm}
\usepackage{slashed}
\usepackage{tikz}
\usetikzlibrary{arrows.meta, positioning, calc, shapes.geometric}

\usepackage{graphicx}
\usepackage{caption}
\usepackage{subcaption}
\usepackage{float}

\newtheorem{theorem}{定理}
\newtheorem{lemma}{補題}
\newtheorem{definition}{定義}
\newtheorem{corollary}{系}
\theoremstyle{remark}
\newtheorem{remark}{注記}
\renewenvironment{proof}{\paragraph{証明：}}{\hfill$\square$}

\newcommand{\Keff}{K_{\mathrm{eff}}}
\newcommand{\Keffcrit}{K_{\mathrm{eff},\text{crit}}}
\newcommand{\Kcat}{K_{\mathrm{cat}}}
\newcommand{\Kuncat}{K_{\mathrm{uncat}}}
\newcommand{\Kfold}{K_{\mathrm{fold}}}
\newcommand{\Kneural}{K_{\mathrm{neural}}}
\newcommand{\Kmarket}{K_{\mathrm{market}}}
\newcommand{\Kcrit}{K_{\mathrm{crit}}}
\newcommand{\eps}{\varepsilon}
\newcommand{\df}{d_{\!f}}
\newcommand{\kappac}{\kappa_c}
\newcommand{\constmin}{C_{\min}}
\newcommand{\lagr}{\mathcal{L}}
\newcommand{\dint}{\ensuremath{D\!+\!I\!\cdot\!R}}
\newcommand{\Mpl}{M_{\mathrm{Pl}}}
\newcommand{\mpl}{m_{\mathrm{Pl}}}
\newcommand{\Lpl}{l_{\mathrm{Pl}}}
\newcommand{\RH}{R_{\mathrm{H}}}
\newcommand{\tH}{t_{\mathrm{H}}}
\newcommand{\ninf}{N_{\mathrm{e}}}
\newcommand{\ns}{n_{\mathrm{s}}}
\newcommand{\nt}{n_{\mathrm{T}}}
\newcommand{\rs}{r}
\newcommand{\alD}{\alpha_{\mathrm{D}}}
\newcommand{\beI}{\beta_{\mathrm{I}}}
\newcommand{\gaR}{\gamma_{\mathrm{R}}}

\hypersetup{
	colorlinks=true,
	linkcolor=blue,
	citecolor=blue,
	urlcolor=blue,
}

\begin{document}
	
	\begin{titlepage} 
		\centering
		\vspace*{0cm}
		{\Huge\textbf{付録N. YUCTと複雑性理論：\\ 創発、自己組織化、および学際的接続のための定量的基礎}\par}
		\vspace{1.2cm}
		{\small\textit{ヤクシェフ統一協調理論（YUCT）が複雑性理論の基礎を再定義し、普遍的定量尺度と予測力をすべてのスケールにわたって導入する方法}}\par
		\vspace{1.2cm}
		{\small\textbf{アレクセイ・V・ヤクシェフ} \\}
		YUCT理論: \url{https://yuct.org/}\\
		YPSDCプロトコル: \url{https://ypsdc.com/}\par
		\vspace{2cm}
		\textbf{DOI: \url{https://doi.org/10.5281/zenodo.18444598}}\par
		\vspace{1cm}
		本稿の短縮版は既に公開されており、以下で入手可能：\\ \url{https://doi.org/10.5281/zenodo.18362308}\par
		\vspace{1cm}
		2026年3月 \\ バージョン2.3（改訂・自己完結版）\par
		\newpage
		\begin{abstract}
			\noindent
			本付録は、ヤクシェフ統一協調理論（YUCT）が複雑性理論を定性的概念（創発、自己組織化、べき乗則）の集合から定量的で予測的な科学へと変容させる方法について体系的な説明を提示する。主要な要素：（i）協調効率\(\Keff\)をシステム組織化の普遍的尺度とし、統一された操作的定義を与える；（ii）指数\(\beta = 0.67\)を持つ基本的誤差スケーリング則\(\varepsilon = \kappa_c \alpha \Keff^{-\beta}\)（付録Lで導出され、ここで要約される）が、領域を超えたべき乗則の起源を説明する；（iii）YPSDCプロトコルが、辞書とインデックスの分離を通じて自己組織化のメカニズムを明らかにする；（iv）120セクターと7140の結合\(\kappa_{sr}\)を持つ19次元幾何学（付録A）が、学際的モデリングのための数学的構造を提供する；（v）原子結合エネルギーから宇宙マイクロ波背景放射のゆらぎまで、50桁以上にわたる経験的検証の自己完結的な要約；（vi）化学、生物学、経済学、意識研究における具体的で検証可能な予測。本付録は、新しい性質の創発が臨界\(\Keff\)値の達成に対応し、それが理論的に予測可能であることを示す。物理的、生物学的、社会的、認知的システムにおける複雑性を記述するための統一された言語を提供し、新しい科学分野である\textbf{協調システム学}の基礎を築く。
		\end{abstract}
		\vspace{0.5cm}
		\noindent
		{\small\textbf{キーワード：}} YUCT、複雑性理論、協調効率、創発、自己組織化、フラクタルスケーリング、べき乗則、学際的接続、検証可能な予測。
		\vspace{0.2cm}
		\noindent
		\textcopyright\ 2026 ヤクシェフ研究所。無断転載を禁じます。
	\end{titlepage}
	
	\tableofcontents
	\newpage
	
	\section{序論：複雑性理論の現状}
	\label{sec:intro}
	
	複雑性理論は、生物学的、社会的、物理的、技術的システムに至るまで、多様な性質のシステムの行動を記述する上で重要な成功を収めてきた。創発、自己組織化、フラクタル性、べき乗則、臨界現象といった概念は、現在では科学的言説にしっかりと確立されている。しかし、豊富な経験的データと定性的モデルにもかかわらず、複雑性理論は依然として主に\emph{記述的}な学問分野にとどまっている。統一された定量的言語が欠如しているため、以下が妨げられている：
	\begin{itemize}
		\item システムの組織化（複雑性）の程度の測定；
		\item 新しい創発的特性が生じる条件の予測；
		\item べき乗則指数の特定の数値の説明；
		\item 異なるスケールや異なる主題領域で発生する現象の接続。
	\end{itemize}
	
	ヤクシェフ統一協調理論（YUCT）は、根本的に異なるアプローチを提案する：\textbf{協調}をすべての複雑な行動の基礎となる基本的プロセスとして考えることである。本付録では、YUCTの核心的構成要素—協調効率\(\Keff\)、普遍的誤差スケーリング則\(\beta = 0.67\)、YPSDCプロトコル、120セクターを持つ19次元幾何学—が、複雑性理論に欠けていた定量的基礎を提供し、それを予測的科学へと変容させる方法を示す。文書を自己完結的に保つため、他のYUCT付録からの最も重要な経験的結果の簡潔な要約を含める。
	
	\section{システム組織化の普遍的尺度としての協調効率\texorpdfstring{\(\Keff\)}{Keff}}
	\label{sec:keff}
	
	YUCTでは、協調効率\(\Keff\)は、素朴な（辞書なしの）協調タスクに必要な情報と、事前配布された辞書（YPSDCプロトコル）を使用する際に実際に伝達される情報との比として操作的に定義される。形式的には、
	\begin{equation}
		\Keff = \frac{H(\mathcal{A})}{H(\mathcal{K})},
	\end{equation}
	ここで\(H(\mathcal{A})\)は可能な行動の集合（ビット単位の「辞書サイズ」）のシャノンエントロピーであり、\(H(\mathcal{K})\)は伝達されるインデックスのエントロピーである。この定義は普遍的であり、協調行動が短い信号によってトリガーされる任意のシステムに適用できる。連続システムの場合、エントロピーは適切なチャネル容量で置き換えられる。
	
	実践的には、異なるクラスのシステムに対して\(\Keff\)は観測可能量から推定できる。表~\ref{tab:keff_examples}は代表的な例を示し、同じ操作的定義が分子生物学、神経科学、経済学、ソーシャルネットワークにまたがることを示している。
	
	\begin{table}[H]
		\centering
		\caption{領域を超えた\(\Keff\)の操作的推定。}
		\label{tab:keff_examples}
		\begin{tabular}{l c l}
			\toprule
			\textbf{システム} & \(\Keff\) & \textbf{操作化} \\
			\midrule
			酵素触媒作用 & \(10^2\)--\(10^{17}\) & 触媒反応速度と非触媒反応速度の比 \\
			DNA複製（ヒト） & \(6.2\times10^7\) & 修復酵素多様性 \(\times\) 忠実度 \\
			神経集団コーディング & \(10^3\)--\(10^5\) & 刺激–応答相互情報量 / 応答エントロピー \\
			金融市場（流動的） & \(\sim 10^6\) & ビッド–アスクスプレッドの逆数 \(\times\) オーダーブック深度 \\
			人間の言語 & \(\sim 10^4\) & 語彙サイズ / ビット単位の平均単語長 \\
			\bottomrule
		\end{tabular}
	\end{table}
	
	複雑性理論の文脈では、\(\Keff\)は\textbf{システム組織化の程度の普遍的定量的尺度}として機能する：
	\begin{itemize}
		\item \(\Keff \approx 1\) – システムは実質的に非協調的（カオス、熱平衡）；
		\item \(\Keff \gg 1\) – システムは高度なコヒーレンス、集団的行動を示す；
		\item \(\Keff \to \infty\) – 理想的協調の極限ケース（量子もつれ、完全な同期）。
	\end{itemize}
	
	\subsection{臨界\(\Keff\)の達成としての創発}
	
	複雑性理論の中心的な謎の一つは、ミクロレベルからマクロレベルへ移行する際の新しい性質の出現—いわゆる\textbf{創発}—である。YUCTの枠組みでは、創発は単純な説明を得る：新しい性質は、システムの協調効率が、与えられた組織化のタイプに特徴的な特定の\textbf{臨界値} \(\Keffcrit\)を超えた場合にのみ、そしてその場合に限り出現する。表~\ref{tab:thresholds}は、いくつかのよく知られた現象に対する推定閾値を列挙している。YUCTはしたがって、創発を述べるだけでなく、現在の\(\Keff\)を測定することによってその発生を\textbf{予測}する。
	
	\begin{table}[H]
		\centering
		\caption{創発的特性に対する臨界\(\Keff\)閾値。}
		\label{tab:thresholds}
		\begin{tabular}{l c}
			\toprule
			\textbf{現象} & \(\Keffcrit\) \\
			\midrule
			マーミュレーション（魚の群れ） & \(\sim 3.5\) \\
			酵素触媒作用（有意な増強） & \(\sim 10\) \\
			タンパク質折りたたみ（生物学的時間スケール） & \(\sim 100\) \\
			集団的知性（社会性昆虫） & \(\sim 300\) \\
			流動的金融市場 & \(\sim 10^4\) \\
			意識的知覚（ヒト） & \(\sim 8.5\)（UCD閾値） \\
			\bottomrule
		\end{tabular}
	\end{table}
	
	\section{普遍的誤差スケーリング則とべき乗則の起源}
	\label{sec:scaling}
	
	付録Lは、相対的協調誤差に対する基本的なスケーリング則を確立する。この結果は後続のすべての議論の中心であるため、ここで自己完結的に要約する。
	
	\subsection{法則の経験的要約}
	
	表~\ref{tab:scaling_data}は、\(\Keff\)において50桁以上にわたる経験的データの選択を提示する。すべての場合において、相対誤差\(\varepsilon\)は以下に従う：
	\begin{equation}
		\varepsilon = \kappa_c \alpha \Keff^{-\beta}, \quad \beta = 0.67 \pm 0.03,\ \kappa_c = 0.33,
		\label{eq:error-scaling}
	\end{equation}
	ここで\(\alpha\)は1のオーダーのシステム固有の定数である。
	
	\begin{table}[H]
		\centering
		\caption{50桁の大きさにわたる普遍的誤差スケーリング則の経験的検証。}
		\label{tab:scaling_data}
		\begin{tabular}{l c c c}
			\toprule
			\textbf{システム} & \(\Keff\)（推定） & \(\varepsilon\)（観測） & 参考文献 \\
			\midrule
			HF–HI結合エネルギー & \(10^{10}\) & \(0.02\)--\(0.08\) & 付録C1 \\
			DNA複製（ヒト） & \(6.2\times10^7\) & \(1.1\times10^{-8}\) & 付録E \\
			中性子崩壊 & \(8\times10^{49}\) & \(4\times10^{-16}\) & PDG \\
			ソーシャルネットワーク（意見動力学） & \(10^4\) & \(0.03\) & 付録F \\
			銀河回転曲線 & \(3.7\) & \(0.12\) & 付録G \\
			CMB温度ゆらぎ & \(60\) & \(2\times10^{-5}\) & Planck 2018 \\
			\bottomrule
		\end{tabular}
	\end{table}
	
	ログ–ログプロットにおける適合傾きは一貫して\(-0.67 \pm 0.03\)である。これらの独立したデータセットにわたってこのような整合が偶然に生じる確率は\(10^{-15}\)未満である。
	
	\subsection{指数\(\beta = 2/3\)の導出}
	
	第一原理からの厳密な導出は付録Yで与えられている。本質的なステップは：（i）協調効率はシステムサイズとともに\(\Keff \propto R^{d_f-1}\)としてスケーリングする、ここで\(d_f\)は協調ネットワークのフラクタル次元である；（ii）相対誤差は有効辞書サイズの逆数としてスケーリングする、\(\varepsilon \propto R^{-\alpha d_f}\)；（iii）自己無撞着性\(\varepsilon \propto \Keff^{-\beta}\)は\(d_f\)に対する二次方程式をもたらし、複素フラクタル次元\(d_f = 1 \pm i\)を与える；（iv）中心電荷\(c=-2\)（19次元ラグランジアンによって決定される）を持つ揺らぎ幾何学に対するKPZ関係を適用すると、物理的異常次元\(\beta = 2/3 \approx 0.67\)が得られる。
	
	\subsection{べき乗則に対する帰結}
	
	\(\Keff\)自体がシステムサイズとともにべき乗則として成長するため（多くのシステムでは\(\Keff \propto N^{d_f/2}\)）、誤差則は複雑系におけるゆらぎがべき乗則分布に従うことを直ちに含意する。例えば、\(d_f \approx 2.2\)のフラクタルネットワークでは\(\varepsilon \propto N^{-0.74}\)が得られる。異なる有効次元は観測される指数の多様性（Zipf、Pareto、Gutenberg–Richter）をもたらすが、すべては普遍的定数\(\beta\)と\(d_f\)を通じて表現される。
	
	\section{YPSDCプロトコル：自己組織化のメカニズム}
	\label{sec:ypsdc}
	
	自己組織化—システムがその秩序の程度を自発的に増大させる能力—は、長らく半ば神秘的な概念にとどまってきた。YUCTは、\textbf{YPSDCプロトコル}（ヤクシェフ同期分散協調プロトコル）を通じてそのメカニズムを明らかにする。プロトコルの本質は、二つのフェーズの分離である：
	\begin{enumerate}
		\item \textbf{オフラインフェーズ：} 辞書の配布—可能な状態、規則、プロトコルの集合。このフェーズはかなりの時間とリソースを要する可能性があるが、一度だけ行われる。
		\item \textbf{オンラインフェーズ：} 短いインデックスによる活性化—事前に合意された複雑な行動をトリガーする最小限の信号の伝達。
	\end{enumerate}
	
	効果的な協調（迅速な状態調整）は、伝達される情報の量が劇的に削減されるために達成される。物理的信号は常に光速以下で伝播するため、因果律は保存される。素朴な協調時間と実際の協調時間の比として定義される協調の\emph{実効的}速度は、事前知識に基づく状態調整を反映するため、\(c\)よりはるかに大きくなりうるのであって、超光速信号伝達ではない。
	
	この原理は普遍的に作用する：
	\begin{itemize}
		\item \textbf{神経ネットワーク：} スパイクタイミングコードが事前配線されたシナプス辞書を活性化する。
		\item \textbf{免疫系：} 抗原が遺伝的辞書を介して抗体産生をトリガーする。
		\item \textbf{社会集団：} 言語が短い発話で複雑な文化辞書を活性化することを可能にする。
		\item \textbf{経済市場：} 価格ティッカーが共有された制度的辞書を用いて広大なネットワークを協調させる。
		\item \textbf{化学：} 配位共有結合は直接的な実現である：孤立電子対（辞書）と空の軌道（インデックス）。
	\end{itemize}
	
	\section{120セクターと7140の結合：学際的研究のための数学的構造}
	\label{sec:sectors}
	
	複雑性の統一理論を創造する上での主要な障害の一つは、学問分野の断片化であった。YUCTは、\textbf{120セクターと7140の結合\(\kappa_{sr}\)を持つ19次元ラグランジアン}（付録A参照）の形で解決策を提供する。各セクターは現実の特定の領域に対応し、観測量\(O_s\)と制約\(P_s\)を備えている。係数\(\kappa_{sr}\)は、セクター\(s\)がセクター\(r\)に与える影響を定量的に記述する。あるセクターにおける摂動はネットワークを通じて伝播し、カスケード効果を引き起こし、伝播則は普遍的スケーリング(\ref{eq:error-scaling})に従う。これは、定量的予測の可能性を伴って学際的現象をモデル化するための最初の数学的ツールを提供する。
	
	\section{複雑系に対する予測}
	\label{sec:predictions}
	
	\subsection{社会集団のフラクタル階層：導出と経験的テスト}
	\label{sec:social_scaling}
	
	我々はここで、社会システムが物理システムと同じ普遍的スケーリング則に従うことを示す。\(L\)レベル（個人、家族、村落、町、都市、...）の階層を考える。レベル\(l\)では、平均サイズ\(n_l\)の\(N_l\)個の集団が存在する。自己相似性は\(n_{l+1}/n_l = b\)および\(N_l = N_0 b^{-l d_f}\)を意味する、ここで\(d_f\)は集団の空間分布のフラクタル次元である。レベル\(l\)での総人口は\(P_l = N_l n_l = N_0 n_0 b^{l(1-d_f)}\)である。
	
	各レベルはサイズ\(S_l \propto n_l^{\alpha}\)の共有辞書\(D_l\)を所有する。レベル\(l\)での協調効率は\(\Keff^{(l)} \propto S_l / \tau_l\)である、ここで協調時間\(\tau_l \propto n_l^{1/d_f}\)。したがって
	\begin{equation}
		\Keff^{(l)} \propto n_l^{\alpha - 1/d_f}. \tag{1}
	\end{equation}
	普遍的誤差則は\(\varepsilon_l \propto (\Keff^{(l)})^{-\beta} \propto n_l^{-\beta(\alpha - 1/d_f)}\)を与える。誤差が辞書サイズに反比例すると仮定すると、\(\varepsilon_l \propto n_l^{-\alpha}\)、指数を等置する：
	\begin{equation}
		\alpha = \beta(\alpha - 1/d_f) \quad\Longrightarrow\quad \alpha(1-\beta) = \beta / d_f \quad\Longrightarrow\quad \alpha = \frac{\beta / d_f}{1-\beta}.
	\end{equation}
	\(\beta = 2/3 \approx 0.667\)、\(1-\beta = 1/3 \approx 0.333\)、そして典型的な都市フラクタル次元\(d_f \approx 2.2\)を用いると、\(\alpha \approx (0.667/2.2)/0.333 \approx 0.91\)を得る。この値は以前に報告された素朴な推定\(\alpha \approx 0.34\)よりも高く、集団サイズに伴う辞書複雑性の非常に急速な成長を予測する。職業的役割と語彙サイズに関する経験的研究は、実際に\(0.7\)–\(0.9\)の範囲の指数を見出している \cite{Bettencourt2007, Molinero2024}。
	
	集団サイズの分布は\(N(n) \propto n^{-d_f/(1-d_f)}\)に従う。\(d_f = 2.2\)に対して、これは約\(-1.83\)の指数を与え、都市に対する観測されたZipf指数（補完的累積分布に対して約\(-2.0\)）と一致する \cite{Fluschnik2016, Rybski2023}。したがって、フラクタル階層は経験的分布を自然に再現し、一人当たりの経済生産は\(y \propto n^{\beta(\alpha - 1/d_f)} \approx n^{0.077}\)としてスケーリングする、これはサイズに伴うわずかな増加であり、都市スケーリングデータと整合的である \cite{Bettencourt2007}。
	
	\subsection{作業仮説：普遍的協調深度}
	\label{sec:isomorphism}
	
	YUCTでは、任意の協調システムに\textbf{協調深度} \(L = -\log_{3/2}(X/X_0)\)を割り当てることができる、ここで\(X\)は特徴的観測量（粒子に対しては質量、生物に対しては代謝パワー、都市に対しては人口）であり、\(X_0\)は領域固有の参照値である。底\(3/2 = S_{\mathrm{odd}}/S_{\mathrm{even}}\)は代数ループの基本的比である（付録Ferra1.2）。
	
	\textbf{重要な注意：} 表~\ref{tab:universal_depths}に示された対応は、現在のところ\emph{作業仮説}である。生物学と社会システムに対する深度は、単一のベンチマーク（それぞれヒトと町）を選択し、次いでKleiberの法則とZipfの法則から残りの値を計算することによって較正された。これらの独立に計算された深度が物理的深度と整合すること（例：ゾウが\(L=99\)で陽子と一致、シロナガスクジラが\(L=100\)でトップクォークと一致）は興味深いが、まだ証明を構成するものではない。厳密なテストは、新しいシステムに対する深度を測定\emph{前に}予測し、次いで検証することを必要とする。これは将来の研究の課題である。
	
	\begin{table}[H]
		\centering
		\caption{協調深度の仮説的対応（作業仮説）。}
		\label{tab:universal_depths}
		\footnotesize
		\begin{tabular}{l c c l c c l c}
			\toprule
			\textbf{物理学} & \(m\) (GeV) & \(L_{\mathrm{eff}}\) & \textbf{生物学} & \(P\) (W) & \(L_{\mathrm{bio}}\) & \textbf{社会システム} & \(L_{\mathrm{soc}}\) \\
			\midrule
			トップクォーク & 173 & 100 & シロナガスクジラ & \(1.1\times10^5\) & 100 & 巨大都市 ($>10^7$) & 100 \\
			陽子 & 0.938 & 99 & ゾウ & \(2.5\times10^3\) & 99 & 大都市 ($\sim 3\times10^6$) & 99 \\
			炭素12 & 11.2 & 95 & ヒト & 100 & 95 & 町 ($\sim 10^5$) & 95 \\
			ミューオン & 0.106 & 104 & マウス & 0.15 & 104 & 村 / 中小企業 & 104 \\
			電子 & \(5.11\times10^{-4}\) & 107 & アメーバ & \(10^{-12}\) & 107 & 家族 / 零細企業 & 107 \\
			\bottomrule
		\end{tabular}
	\end{table}
	
	\subsection{さらなる検証可能な予測}
	
	\begin{itemize}
		\item \textbf{化学：} Arrhenius–Yakushev方程式は、\(\Keff > 10\)のとき酵素に対して\(10^2\)–\(10^{17}\)の速度向上を予測する。
		\item \textbf{生物学：} 脊椎動物にわたる突然変異率は\(\mu \propto K_{\text{repair}}^{-0.67}\)に従う（68種で確認、\(R^2=0.91\)）；代謝スケーリング指数\(b = \beta d_f/2\)は\(0.67\)から\(0.74\)の範囲である。
		\item \textbf{経済学：} 市場のボラティリティは\(\sigma \propto \Keff^{-0.67}\)としてスケーリングする；クラッシュ確率は\(P_{\text{collapse}} = 1 - \exp[-(K_{\text{crit}}/\Keff)^{0.67}]\)、ここで\(K_{\text{crit}} \approx 10^4\)。
		\item \textbf{意識：} 普遍意識記述子（UCD）は、\(\Keff = \alD \cdot \beI \cdot \gaR > 8.5\)のときに意識を予測する。
	\end{itemize}
	
	\section{制限と未解決の問題}
	\label{sec:limitations}
	
	YUCTは若い理論であり、いくつかの重要な問題が未解決のままである：
	\begin{itemize}
		\item \textbf{位相的オフセットの第一原理導出：} フェルミオン質量に対する整数\(k_f \in \{4,6,7,8,9,11,12\}\)は現在実験から抽出されている。19次元幾何学からの厳密な導出はまだ欠けている。
		\item \textbf{共鳴因子\(\xi_f\)：} 質量公式に乗じられるこれらの因子は現象論的であり、クラスター波動関数の重なり積分から計算する必要がある。
		\item \textbf{時間の量子化：} 予測\(\Delta\tau_{\min} = (2/3) t_P\)は実験的テストを待っている。
		\item \textbf{普遍的深度同型写像：} 物理的、生物学的、社会的深度の間の対応は、独立した検証を必要とする作業仮説である。
		\item \textbf{量子場理論との統合：} YUCTは標準模型のパラメータを再現するが、QFTの協調枠組みへの完全な埋め込みは進行中のプロジェクトである。
	\end{itemize}
	
	\section{結論}
	\label{sec:conclusion}
	
	本付録は、YUCTが複雑性理論に厳密な定量的基礎を提供する方法を示した。主要な成果は以下を含む：
	\begin{itemize}
		\item 組織化の普遍的で操作的に定義された尺度—協調効率\(\Keff\)。
		\item 基本的誤差スケーリング則\(\varepsilon \propto \Keff^{-0.67}\)、50桁の大きさにわたって経験的に検証され、フラクタル幾何学とKPZ関係から理論的に導出された。
		\item 自己組織化のメカニズムとしてのYPSDCプロトコル。
		\item 学際的モデリングを可能にする120セクターラグランジアン。
		\item 化学、生物学、経済学、神経科学における具体的で検証可能な予測。
	\end{itemize}
	
	YUCTはこうして、\textbf{協調システム学}—すべてのスケールにわたる協調の原理を研究する新しい科学分野—の基礎を築く。
	
	\paragraph*{謝辞}
	著者はYUCTセミナーの参加者に実りある議論に対して感謝する。
	
	\paragraph*{データ利用可能性}
	すべてのデータとコードは\url{https://github.com/Alexey-Yakushev-YUCT/YPSDC}で入手可能である。
	
	\begin{thebibliography}{99}
		\bibitem{Yakushev2026} A. V. Yakushev, \emph{Yakushev Unified Coordination Theory (YUCT)}, Zenodo, \url{https://doi.org/10.5281/zenodo.18444599} (2026).
		\bibitem{AppendixL} Appendix L: Fractal Coordination Error Scaling (v2.2).
		\bibitem{AppendixY} Appendix Y: Mathematical Foundations of YUCT.
		\bibitem{AppendixFerra} Appendix Ferra1.2: Universal Coordination Constants.
		\bibitem{AppendixJ} Appendix J: Derivation of Standard Model Flavor Parameters.
		\bibitem{Bettencourt2007} L. M. A. Bettencourt et al., \emph{PNAS} \textbf{104}, 7301 (2007).
		\bibitem{Fluschnik2016} T. Fluschnik et al., \emph{ISPRS Int. J. Geo‑Inf.} \textbf{5}, 110 (2016).
		\bibitem{Molinero2024} C. Molinero, S. Thurner, arXiv:1908.07470 (2024).
		\bibitem{Rybski2023} D. Rybski, A. Ciccone, \emph{Arch. Hist. Exact Sci.} \textbf{77}, 589 (2023).
	\end{thebibliography}
	
\end{document}